\section{Introducción}
En una bodega con estanterías metálicas, paneles y muros de concreto, la señal de un nodo sensor no llega al receptor por un único camino: a la componente directa se suman réplicas reflejadas, atenuadas y retardadas. Las mediciones clásicas de canales interiores describen este fenómeno como llegadas en grupos con dispersión de retardo eficaz de decenas a cientos de nanosegundos \cite{saleh1987,hashemi1993}; en plantas industriales se han medido valores entre 30 y 300~ns, influidos por el inventario y los materiales de construcción \cite{rappaport1989}. Cuando el retardo de la réplica es comparable con la duración del símbolo, el eco se mezcla con el símbolo siguiente y aparece interferencia intersimbólica (ISI); en frecuencia, el canal deja de ser plano y produce desvanecimiento selectivo \cite{proakis2008,rappaport2002}.

El Grupo~5 eligió el escenario S3 de la guía del curso \cite{guia1}: un sensor transmite variables de baja velocidad (estado, temperatura, vibración) hacia una estación de supervisión dentro de una bodega, y la reflexión dominante deforma los símbolos. El objetivo es aislar ese efecto, sin otras fuentes de ISI, para observarlo y medirlo de forma controlada. Este documento presenta el escenario y la decisión de diseño (Sección~III), la arquitectura implementada en GNU Radio, que incluye el receptor exigido en la Guía~2 \cite{guia2}, las cuatro pruebas de perturbación y la caracterización BER de la Fase~4 \cite{guia4} (Sección~IV), y cierra con las conclusiones.

\textbf{Ajuste a la ficha de decisión.} En la ficha ``Decisión Técnica Ciega'' del grupo, el escenario (S3), la perturbación (camino directo + eco), la variable (ganancia compleja del eco) y la justificación escrita corresponden a multitrayectoria. Sin embargo, quedaron marcados el problema P1 y la hipótesis H1, que corresponden a perturbaciones impulsivas. Para que el informe sea coherente se adoptó explícitamente la hipótesis equivalente para multitrayectoria: \emph{al aumentar $|\alpha|$, la constelación recuperada se desdobla, el ojo se cierra en proporción $1-|\alpha|$ y la BER crece para un mismo $E_b/N_0$}.
